\section{Related work and contribution boundary}

\paragraph{Direct preference optimization.}
DPO converts a KL-regularized preference-learning problem into a logistic
classification objective over policy and reference log ratios
\citep{rafailov2023dpo}. Our work does not propose a new optimizer. It studies a
special case in which both alternatives are single verdict tokens and the
preference is determined by an existing binary label.

\paragraph{Preference training for guards.}
GuardReasoner combines reasoning SFT with hard-sample DPO
\citep{liu2025guardreasoner}. AIMS studies intent-aware supervision across SFT,
DPO, distillation, and reinforcement learning \citep{ferrao2026aims}. DT-Guard
uses rollout-guided hard-case optimization in a reasoning-active,
reasoning-free-inference guard \citep{liu2026dtguard}. These systems use richer
reasoning, intent, or rollout signals and pursue absolute guard quality. Our
question is orthogonal: with information held constant, is reference centering
better than the corresponding uncentered pairwise objective?

\paragraph{Where this sits: reference centering as an anti-forgetting dial.}
The reduction of \Cref{lem:rescale} places this study in a family it is usually
not discussed with. A frozen reference term penalises movement away from the
starting policy, which is what an explicit KL penalty, elastic-weight-style
regularisation, and an output-space average of base and adapter also do by
different means. Read that way, the question is not ``is DPO a better guard
objective'' but ``does an \emph{implicit} anchor purchase the same
represented/transfer trade the \emph{explicit} anchors purchase, and at what
price?'' This matters for interpretation in both directions. If
$F_{\mathrm{ref}}$ is positive, reference centering is a zero-extra-inference
anti-forgetting mechanism, cheaper at serving time than averaging two models. If
it is a bounded null, then a pipeline that reaches for DPO on binary guard labels
is paying preference-optimization complexity for a gradient reweighting, and an
explicit anchor with a tunable strength is the better-instrumented choice. Either
outcome is a usable finding, which is why the design fixes how a null is reported
before running.

\paragraph{Guard evaluation.}
We retain a multi-source evaluation spanning real-world toxicity, jailbreak,
over-refusal, and safety stress suites \citep{lin2023toxicchat,chao2024jailbreakbench,
rottger2024xstest,han2024wildguard,cui2024orbench,mazeika2024harmbench}. The
benchmark panel is used retrospectively; a future confirmatory claim requires a
separately sealed cohort. Because the parent suite also supplies the seed variance
that sets this design's resolving power (\Cref{tab:power}), the same reuse that
makes the evaluation retrospective is what makes the power calculation possible
before any Paper C run exists.

